% Copia mecánica íntegra de apendices.tex:1--761.
% Residencia material del ensamblaje sucesor de 102 capítulos.
\part{Apéndices demostrativos y tablas de clasificación}

\chapter{Cálculo afín ordenado de las hojas APP}

\section{Semiproducto afín y composición de rutas}

Sea \(U_{\APP}\) el módulo radial anterior a su evaluación
arquimediana. La envolvente correcta es el monoide
\[
 \AffEnd(U_{\APP})
 =\operatorname{End}(U_{\APP})\ltimes U_{\APP},
\]
pues las hojas de umbral o plegamiento pueden contener endomorfismos no
invertibles. La acción de \((\Lambda,C)\) sobre \(u\in U_{\APP}\) es
\[
 (\Lambda,C)\cdot u=\Lambda u+C.
\]
La composición se calcula sin conmutar factores:
\begin{align}
 (\Lambda_2,C_2)(\Lambda_1,C_1)\cdot u
 &=\Lambda_2(\Lambda_1u+C_1)+C_2 \notag\\
 &=(\Lambda_2\Lambda_1)u+(C_2+\Lambda_2C_1).
 \label{espiral:eq:ap-comp-afin}
\end{align}
Por tanto,
\[
 (\Lambda_2,C_2)(\Lambda_1,C_1)
 =(\Lambda_2\Lambda_1,C_2+\Lambda_2C_1).
\]
La segunda coordenada no es una suma ordinaria de desplazamientos: es un
cociclo torcido por la acción lineal acumulada.

\begin{proposition}[Fórmula de los prefijos]
Para una ruta de \(N\) factores \(h_j=(\lambda_j,C_j)\), ordenados desde
\(j=0\) hasta \(j=N-1\), la holonomía prefija es
\[
 H_N=h_{N-1}\cdots h_1h_0=(\Lambda_N,\mathcal C_N),
\]
donde
\[
 \Lambda_N=\prod_{j=0}^{N-1}\lambda_j
\]
y
\[
 \mathcal C_N
 =C_{N-1}+\lambda_{N-1}C_{N-2}
 +\lambda_{N-1}\lambda_{N-2}C_{N-3}
 +\cdots+
 \lambda_{N-1}\cdots\lambda_1C_0.
\]
\end{proposition}
\begin{proof}
Para \(N=1\) la fórmula es inmediata. Si vale para \(N\), entonces
\[
 h_NH_N
 =(\lambda_N\Lambda_N,C_N+\lambda_N\mathcal C_N),
\]
que produce exactamente la expresión de longitud \(N+1\). La inducción
preserva el orden de la ruta y no identifica recorridos con los mismos
extremos.
\end{proof}

\section{Diferencia suma--producto como curvatura de orden}

Sean \(A_a=(I,a)\) y \(M_\lambda=(\lambda,0)\). En el sector invertible,
\[
 A_aM_\lambda A_a^{-1}M_\lambda^{-1}
 =A_{a(1-\lambda)}.
\]
La demostración directa sobre \(u\) es
\[
 u\xmapsto{M_\lambda^{-1}}\lambda^{-1}u
 \xmapsto{A_a^{-1}}\lambda^{-1}u-a
 \xmapsto{M_\lambda}u-\lambda a
 \xmapsto{A_a}u+a(1-\lambda).
\]
Así, la diferencia entre acumular y plegar queda publicada como una
traslación. Cuando \(\lambda=1\), el conmutador se anula; cuando
\(\lambda\neq1\), la ruta conserva una holonomía que no puede reconstruirse
a partir del punto final solamente.

\begin{example}[Dos rutas con el mismo dato visible]
Fijemos \(u_0=0\), \(a=1\), \(\lambda=2\). Las rutas
\[
 A_1M_2:u_0\longmapsto1,
 \qquad
 M_2A_{1/2}:u_0\longmapsto1
\]
publican el mismo valor final. Sin embargo, sus ledgers contienen factores,
desplazamientos y prefijos distintos. La evaluación puntual identifica las
rutas; la holonomía afín enriquecida las separa.
\end{example}

\section{Tabla exhaustiva de la carta local de 2 orientaciones}

La tabla siguiente no postula que todo grado radial global proceda de un
solo par. Registra la carta elemental obtenida al componer dos orientaciones
TRIT. El grado de una ruta general es el cociclo entero de
\cref{espiral:eq:cociclos-radiales}; esta carta proporciona sus nueve generadores
locales.

\begin{longtable}{rrrrrrp{.27\textwidth}}
\caption{Los \(9\) pares orientados de la carta TRIT local.}
\label{espiral:tab:ap-nueve-pares}\\
\toprule
\(\tau_+\)&\(\tau_-\)&\(p\)&\(a\)&\(r_3\)&\(c_3\)&Lectura local\\
\midrule
\endfirsthead
\toprule
\(\tau_+\)&\(\tau_-\)&\(p\)&\(a\)&\(r_3\)&\(c_3\)&Lectura local\\
\midrule
\endhead
\( 1\)&\( 1\)&\( 2\)&\( 0\)&\(-1\)&\( 1\)&apertura radial doble con carry positivo\\
\( 1\)&\( 0\)&\( 1\)&\( 1\)&\( 1\)&\( 0\)&apertura y umbral\\
\( 1\)&\(-1\)&\( 0\)&\( 2\)&\( 0\)&\( 0\)&equilibrio con orientación positiva\\
\( 0\)&\( 1\)&\( 1\)&\(-1\)&\( 1\)&\( 0\)&umbral y apertura\\
\( 0\)&\( 0\)&\( 0\)&\( 0\)&\( 0\)&\( 0\)&umbral doble\\
\( 0\)&\(-1\)&\(-1\)&\( 1\)&\(-1\)&\( 0\)&umbral y contracción\\
\(-1\)&\( 1\)&\( 0\)&\(-2\)&\( 0\)&\( 0\)&equilibrio con orientación negativa\\
\(-1\)&\( 0\)&\(-1\)&\(-1\)&\(-1\)&\( 0\)&contracción y umbral\\
\(-1\)&\(-1\)&\(-2\)&\( 0\)&\( 1\)&\(-1\)&contracción doble con carry negativo\\
\bottomrule
\end{longtable}

Aquí \(p=r_3+3c_3\), con residuo balanceado
\(r_3\in\{-1,0,1\}\). El carry distingue \(p=2\) de su residuo \(-1\) y
\(p=-2\) de su residuo \(1\). El par
\((p,a)\) distingue asimismo las tres genealogías locales de
grado \(0\).

\chapter{Certificado nonádico y geometría de la memoria}

\section{Distinción tipada entre catálogo decimal, palabras ternarias y fibra ambiente}

\begin{longtable}{p{.18\textwidth}p{.18\textwidth}p{.52\textwidth}}
\caption{Estratos del certificado de la conexión nonádica conjunta.}
\label{espiral:tab:certificado-nonadico}\\
\toprule
Objeto&Cardinalidad&Función matemática\\
\midrule
\endfirsthead
\toprule
Objeto&Cardinalidad&Función matemática\\
\midrule
\endhead
Catálogo de ternas de origen&\(104\,976\)&Dominio finito previo a la
selección conjunta; no es una tabla decimal.\\
Emisiones decimales \(U_6\)&\(468\)&Emisiones compatibles de longitud \(6\)
en la publicación decimal.\\
Palabras visibles \(w_6\)&\(243\)&Palabras ternarias publicadas tras el
tipado; no son el ambiente completo.\\
Ambiente \(\F_3^6\)&\(729\)&Espacio de todas las palabras ternarias de
longitud \(6\); contiene a las \(243\) visibles.\\
Niveles certificados&\(396\)&Prolongación finita hasta las fronteras
declaradas.\\
Trits recorridos&\(2\,376\)&Longitud total \(396\cdot6\).\\
Vueltas completas&\(43\)&Ciclos nonádicos completos en el tramo
certificado.\\
Pares \((K,K+9)\)&\(387\)&Pares de igual fase por canal; el ledger puede
diferir.\\
\bottomrule
\end{longtable}

La cadena de prolongación se conserva en el orden
\[
 w_6\longrightarrow w_{12}\longrightarrow w_{18}
 \longrightarrow w_{24}\longrightarrow w_{30}
 \longrightarrow R_{36}\longrightarrow G_9.
\]
Ninguno de estos objetos puede sustituirse por el cociente modal posterior.
La holonomía \(\Gamma_9\) actúa sobre la dinámica enriquecida; la
publicación de memoria reducida conserva sólo la información especificada
por su codominio.

\section{Retorno de fase, avance de memoria y estado completo}

Con
\[
 \rho_9(N)=1+((N-1)\bmod9),
 \qquad
 q_9(N)=\left\lfloor\frac{N-1}{9}\right\rfloor,
\]
se tiene
\[
 \rho_9(N+9)=\rho_9(N),
 \qquad q_9(N+9)=q_9(N)+1.
\]
Por ello,
\[
 H_9(N+9)=H_9(N)+(0,0,1).
\]
La igualdad de la primera coordenada de fase no implica igualdad del estado:
\[
 g(K+9)=g(K),
 \qquad B_{K+9}\neq B_K
\]
en general. La fórmula helicoidal es precisamente la representación mínima
de esta distinción.

\section{Extensión dodecafásica no escindida}

La sucesión
\[
 0\longrightarrow C_{12}\longrightarrow C_{108}
 \longrightarrow C_9\longrightarrow0
\]
acopla el calendario dodecafásico con el retorno nonádico. Si existiera una
sección homomorfa global \(s:C_9\to C_{108}\), el generador de \(C_9\)
debería elevarse a un elemento de orden divisor de \(9\) que proyectase en el
generador. El lift temporal activo conserva, sin embargo, el desplazamiento
dodecafásico acumulado y no identifica el estado tras una vuelta de fase.
La costura se registra como carry; geométricamente se manifiesta como
traslación axial de la hélice.

\section{5 regiones de \texorpdfstring{\(\pi\)}{pi} y no fidelidad de la evaluación}

La microfibra de \(\pi\) contiene cinco regiones orientadas con
multiplicidad total
\[
 432+4\cdot144=1008.
\]
Las regiones publican el mismo valor real, pero conservan genealogías
distintas. Este resultado proporciona el precedente interno de la
clasificación espiral: una curva visible no determina por sí sola la región,
la hoja, el carry ni la memoria que la produjo.

\chapter{Álgebra potencia--logarítmica y 5 publicaciones canónicas}

\section[Coordenada potencia--logarítmica e inversa]{La coordenada \(L_p\) y su inversa}

Para \(x>0\), se define
\[
 L_p(x)=
 \begin{cases}
 (x^p-1)/p,&p\neq0,\\
 \log x,&p=0.
 \end{cases}
\]
Su inversa es
\[
 \exp_p(u)=
 \begin{cases}
 (1+pu)^{1/p},&p\neq0,\ 1+pu>0,\\
 \e^u,&p=0.
 \end{cases}
\]
La derivada
\[
 \frac{\dd}{\dd x}L_p(x)=x^{p-1}>0
\]
prueba que \(L_p\) es estrictamente creciente en su dominio y, por tanto,
invertible.

\begin{proposition}[Ley suma--producto]
Para \(x,y>0\),
\[
 L_p(xy)=L_p(x)+L_p(y)+pL_p(x)L_p(y).
\]
\end{proposition}
\begin{proof}
Si \(p\neq0\),
\begin{align*}
 L_p(x)+L_p(y)+pL_p(x)L_p(y)
 &=\frac{x^p-1}{p}+\frac{y^p-1}{p}
   +\frac{(x^p-1)(y^p-1)}p\\
 &=\frac{x^py^p-1}{p}=L_p(xy).
\end{align*}
Para \(p=0\), la igualdad se reduce a
\(\log(xy)=\log x+\log y\).
\end{proof}

\begin{corollary}[Grupo deformado en el dominio positivo]
La operación
\[
 u\boxplus_pv=u+v+puv
\]
satisface
\[
 1+p(u\boxplus_pv)=(1+pu)(1+pv).
\]
Su neutro es \(0\) y el inverso de \(u\) es
\[
 \ominus_pu=-\frac{u}{1+pu},
\]
si \(1+pu\neq0\).
\end{corollary}

\section{Derivación de las 5 curvas canónicas}

Sea \(m\) la profundidad publicada y supongamos
\[
 L_p(r/r_0)=\beta m,
 \qquad
 \theta=\theta_0+\omega m,
 \qquad\omega\neq0.
\]
Eliminando \(m=(\theta-\theta_0)/\omega\),
\[
 r(\theta)=r_0
 \left[1+p\frac{\beta}{\omega}
 (\theta-\theta_0)\right]^{1/p},
 \qquad p\neq0,
\]
y
\[
 r(\theta)=r_0
 \exp\!\left[\frac{\beta}{\omega}
 (\theta-\theta_0)\right]
\]
para \(p=0\). Las reoriginaciones y reescalados positivos producen las
formas convencionales:

\begin{longtable}{rllp{.31\textwidth}}
\caption{Publicaciones canónicas del sector traslacional.}
\label{espiral:tab:cinco-publicaciones-detalle}\\
\toprule
\(p\)&Ecuación normal&Nombre convencional&Dominio y rasgo\\
\midrule
\endfirsthead
\toprule
\(p\)&Ecuación normal&Nombre convencional&Dominio y rasgo\\
\midrule
\endhead
\(2\)&\(r^2=a^2\theta\)&Fermat&\(\theta\ge0\); incremento lineal del área
radial cuadrática.\\
\(1\)&\(r=a\theta\)&Arquímedes&\(\theta\ge0\); incremento radial
constante por unidad angular.\\
\(0\)&\(r=r_0\e^{b\theta}\)&logarítmica&\(\theta\in\R\); dilatación
constante por unidad angular.\\
\(-1\)&\(r=a/\theta\)&hiperbólica&\(\theta>0\); producto \(r\theta\)
constante.\\
\(-2\)&\(r^2=a^2/\theta\)&lituus&\(\theta>0\); producto
\(r^2\theta\) constante.\\
\bottomrule
\end{longtable}

\section{Inversión de hojas}

Para \(p\neq0\),
\[
 L_{-p}(x^{-1})
 =\frac{x^p-1}{-p}=-L_p(x).
\]
La identidad también vale para \(p=0\). Por ello, la inversión radial
intercambia \(p\) con \(-p\) y cambia el signo de la coordenada
linealizada. Esta relación agrupa Fermat con el lituus y Arquímedes con la
hiperbólica como orientaciones conjugadas del lector, sin identificarlas
como curvas iguales.

\section{Sector afín general}

Si la coordenada radial satisface
\[
 u_{n+1}=\Lambda u_n+C,
\]
la iteración interna exacta es
\[
 u_n=\Lambda^n u_0+\sum_{k=0}^{n-1}\Lambda^kC.
\]
Si \(I-\Lambda\) es invertible,
\(u_*=(I-\Lambda)^{-1}C\) y
\(u_n-u_*=\Lambda^n(u_0-u_*)\). Sólo en una realización escalar positiva,
\(\chi\Lambda=\lambda\chi\) con \(\lambda>0\), se define
\(u_*=C/(1-\lambda)\). Con \(\theta_n=\theta_0+n\omega\),
\(\omega\neq0\), la interpolación publicada en esa realización es
\[
 r(\theta)=r_0\exp_p\!\left[
 u_*+\lambda^{(\theta-\theta_0)/\omega}(u_0-u_*)
 \right].
\]
La taxonomía completa no es, pues, un listado de cinco nombres. Contiene al
menos el grado \(p\), la orientación \(a\), la clase afín
\([\Lambda,C]\), el cociclo de carry, la frontera, el incremento angular y
el ledger de memoria.

\chapter{Naturalidad, equivariancia y no fidelidad}

\section{El lector de prefijos}

Para una historia truncada
\[
 \widehat\gamma_N=\widehat e_N\cdots\widehat e_1,
\]
el lector enriquecido registra
\[
 \mathcal P_{\rad,N}^{\enr}(\widehat\gamma_N)
 =
 \left(
 (H_0,\Led_0);
 \bigl(p(\widehat e_j),a(\widehat e_j),P_j,A_j,
 h_{e_j},H_j,\Led_j\bigr)_{j=1}^{N}
 \right).
\]
El término \(H_j=h_{e_j}\cdots h_{e_1}\) conserva los prefijos. El ledger
\(\Led_j\) conserva, según la ruta, hoja, residuo, cociente, carry,
frontera, orientación, cilindro, firma, supervivencia y memoria. Las
coordenadas reales se obtienen después:
\[
 \mathcal P_{\arq}^{\mathrm{esp}}
 =\ev_{\R^2}\circ\mathcal P_{\rad}^{\enr}.
\]

\section{Prueba completa de naturalidad bajo truncamiento}

Sea \(\pi_{N,M}\) el borrado del sufijo de longitud \(N-M\), con \(M<N\).
Sea \(\operatorname{pr}_{N,M}\) la proyección que conserva el registro
inicial y los primeros \(M\) registros de arista. Entonces
\[
 \operatorname{pr}_{N,M}\,
 \mathcal P_{\rad,N}^{\enr}(\widehat\gamma_N)
 =
 \left(
 (H_0,\Led_0);
 \bigl(p(\widehat e_j),a(\widehat e_j),P_j,A_j,
 h_{e_j},H_j,\Led_j\bigr)_{j=1}^{M}
 \right).
\]
Por definición de truncamiento,
\[
 \pi_{N,M}(\widehat\gamma_N)=\widehat e_M\cdots\widehat e_1,
\]
y al volver a aplicar el lector se obtiene la misma secuencia. Por tanto,
\[
 \operatorname{pr}_{N,M}\circ
 \mathcal P_{\rad,N}^{\enr}
 =
 \mathcal P_{\rad,M}^{\enr}\circ\pi_{N,M}.
\]
La demostración no usa el valor de la curva. Usa la identidad de prefijos,
que es la información que el lector reducido de siete coordenadas habría
perdido.

\section{Refinamientos no triviales}

Para un refinamiento \(r:\gamma\to\gamma'\), la naturalidad requiere un
mapa \(\widehat r\) sobre ledgers y la igualdad local
\[
 h_e
 =
 h_{e'_k}\cdots h_{e'_1}
\]
para cada arista \(e\) sustituida por la cadena refinada
\((e'_1,\ldots,e'_k)\). Si además el carry, la frontera y la memoria de la
cadena refinada componen los registros del paso original, entonces
\[
 \widehat r\circ\mathcal P_{\rad}^{\enr}
 =
 \mathcal P_{\rad}^{\enr}\circ r.
\]
Esta condición es un cuadrado verificable; no debe suponerse para un
refinamiento arbitrario sin exhibir los factores.

\section{Equivarianza, no invariancia}

La acción de \(\Gamma_9\) devuelve la fase y prolonga el estado. En el
codominio del lector actúa
\(\widehat\Gamma_9^{\rad}\), que desplaza los registros y actualiza el
último prefijo. La identidad correcta es
\[
 \mathcal P_{\rad}^{\enr}(\Gamma_9\gamma)
 =
 \widehat\Gamma_9^{\rad}
 \mathcal P_{\rad}^{\enr}(\gamma).
\]
Exigir invariancia borraría precisamente la memoria que la conexión
conserva. En el sector reversible, la inversión de rutas satisface
\[
 H_{C_{\mathrm{rev}}\gamma}=H_\gamma^{-1}
\]
cuando todos los factores pertenecen al subgrupo afín invertible. Esta
afirmación tipada no se extiende automáticamente a una involución global
sobre cualquier cociente posterior.

\section{3 demostraciones de no fidelidad}

\begin{proposition}[Circunferencia intrínseca y circunferencia proyectada]
Existen dos historias con la misma curva visible
\[
 c(\theta)=(\cos\theta,\sin\theta)
\]
y ledgers distintos.
\end{proposition}
\begin{proof}
La primera historia posee memoria axial constante \(m=0\). La segunda es
la hélice
\[
 h(\theta)=(\cos\theta,\sin\theta,\theta/2\pi),
\]
cuyo ledger registra \(m(\theta)=\theta/2\pi\). Ambas satisfacen
\(\pi_{xy}\circ h=c\). La proyección es, por tanto, no inyectiva.
\end{proof}

\begin{proposition}[Mismo extremo, distinto orden afín]
Si \(a(1-\lambda)\neq0\), los caminos \(A_aM_\lambda\) y
\(M_\lambda A_a\) no son equivalentes en el lector enriquecido, aunque
puedan calibrarse para publicar un mismo extremo.
\end{proposition}
\begin{proof}
Su conmutador es \(A_{a(1-\lambda)}\neq I\). El dato del conmutador
pertenece al ledger de prefijos y desaparece al conservar sólo el extremo.
\end{proof}

\begin{proposition}[Misma curva, distinta región genealógica]
La igualdad del valor real publicado no identifica las cinco regiones de
la microfibra de \(\pi\); por consiguiente, tampoco basta para identificar
la clase espiral cuando el lector conserva región, frontera o carry.
\end{proposition}

\chapter{Taxonomía estratificada de correspondencias espirales}

\section{Clasificación por nivel matemático}

\begin{longtable}{p{.20\textwidth}p{.30\textwidth}p{.39\textwidth}}
\caption{Niveles que no deben colapsarse en una sola noción de curvatura.}
\label{espiral:tab:niveles-curvatura}\\
\toprule
Nivel&Invariante&Significado\\
\midrule
\endfirsthead
\toprule
Nivel&Invariante&Significado\\
\midrule
\endhead
APP&\(F(S,P)=1,\ F(P,S)=-1\)&Curvatura discreta del orden de acumulación
y plegamiento.\\
TRIT&\(\tau\in\{-1,0,1\}\), \(J_\tau^2=-\tau I\)&Régimen local elíptico,
parabólico o hiperbólico del transporte.\\
TPK&\(\Gamma_9\), carry, frontera, ruta y ledger&Holonomía de fase y
memoria sobre el estado enriquecido.\\
Radial&\(p,a,[\Lambda,C]\)&Caracteres que determinan la ley
radial antes de reconocer una curva.\\
Arquimediano&\(r(\theta)\)&Curva publicada en \(\R^2\); puede olvidar
genealogía.\\
Frenet&\(\kappa_{\mathrm{F}},\mathcal T_{\mathrm{F}}\)&Curvatura y torsión
diferenciales de la curva o suspensión visible.\\
Dimensional&campos, escalas y representaciones&Realización física posterior,
con premisas propias.\\
\bottomrule
\end{longtable}

\section{Clasificación por giro, radio y memoria}

\begin{longtable}{ccccl}
\caption{Taxonomía cinemática completa de los casos degenerados mínimos.}
\label{espiral:tab:taxonomia-cinematica-ap}\\
\toprule
\(\Delta\theta\)&\(\Delta u_{\rad}\)&\(\Delta m\)&Salida visible&
Información genealógica\\
\midrule
\endfirsthead
\toprule
\(\Delta\theta\)&\(\Delta u_{\rad}\)&\(\Delta m\)&Salida visible&
Información genealógica\\
\midrule
\endhead
\(0\)&\(0\)&\(0\)&punto&estado estacionario publicado\\
\(0\)&\(\neq0\)&\(0\)&recta radial&cambio de escala\\
\(0\)&\(0\)&\(\neq0\)&fibra axial&memoria sin desplazamiento transversal\\
\(0\)&\(\neq0\)&\(\neq0\)&hoja radial suspendida&escala y memoria\\
\(\neq0\)&\(0\)&\(0\)&circunferencia&cierre angular intrínseco\\
\(\neq0\)&\(0\)&\(\neq0\)&hélice cilíndrica&fase cerrada, memoria elevada\\
\(\neq0\)&\(\neq0\)&\(0\)&espiral plana&carácter radial visible\\
\(\neq0\)&\(\neq0\)&\(\neq0\)&suspensión helicoidal radial&giro, radio y
profundidad conservados\\
\bottomrule
\end{longtable}

\section{Clasificación por holonomía afín}

La tabla siguiente pertenece a una realización escalar real de la parte
lineal, \(\chi\Lambda=\lambda\chi\); para aligerar su lectura se escribe
\(\Lambda\) en la columna de condiciones. No clasifica endomorfismos
abstractos por desigualdades numéricas.

\begin{longtable}{p{.18\textwidth}p{.22\textwidth}p{.48\textwidth}}
\caption{Tipos de holonomía sobre la coordenada radial.}
\label{espiral:tab:tipos-holonomia}\\
\toprule
Condición&Órbita de \(u\)&Interpretación\\
\midrule
\endfirsthead
\toprule
Condición&Órbita de \(u\)&Interpretación\\
\midrule
\endhead
\(\Lambda=1,\ C=0\)&\(u_n=u_0\)&radio constante; la memoria puede ser
nula u oculta por proyección.\\
\(\Lambda=1,\ C\neq0\)&\(u_n=u_0+nC\)&sector traslacional que publica la
familia \(L_p\).\\
\(0<\Lambda<1\)&\(u_n\to u_*\)&contracción hacia una sección fija.\\
\(\Lambda>1\)&\(|u_n-u_*|\) crece&expansión desde una sección fija.\\
\(\Lambda<0\)&alternancia de orientación&requiere controlar el dominio de
\(\exp_p\) en cada paso.\\
\(\Lambda=0\)&\(u_{n+1}=C\)&plegamiento no invertible; pertenece al
monoide, no al grupo afín.\\
\bottomrule
\end{longtable}

\section{Criterio de equivalencia genealógica}

Dos rutas no se identifican por publicar la misma curva. Se exige un
isomorfismo tipado que preserve término a término:
\[
 \bigl(
 (x_0,H_0,\Led_0);
 (\widehat e_j,p(\widehat e_j),a(\widehat e_j),P_j,A_j,
 h_{e_j},H_j,\Led_j)_{j=1}^{N}
 \bigr)
\]
y que conmute con truncamientos y refinamientos admitidos. La clasificación
resultante es más fina que la taxonomía convencional porque incorpora la
historia que precede a la aparición de \(\R^2\).

\chapter{Realizaciones posteriores y controles de tipo}

\section{Modo electromagnético circular}

Sea
\[
 \zeta=\frac{z-c_{\mathrm{vac}}t}{\lambda},
\]
y
\[
 \mathbf E_\sigma
 =E_0(\mathbf e_1\cos2\pi\zeta
 +\sigma\mathbf e_2\sin2\pi\zeta),
 \qquad
 \mathbf B_\sigma
 =c_{\mathrm{vac}}^{-1}\widehat{\mathbf k}\times\mathbf E_\sigma.
\]
Entonces
\[
 \mathbf E_\sigma\perp\mathbf B_\sigma\perp\widehat{\mathbf k},
 \quad |E|=c_{\mathrm{vac}}|B|,
 \quad\omega=c_{\mathrm{vac}}k.
\]
A tiempo fijo, el extremo normalizado del campo en el fibrado
fase--posición describe
\[
 (\cos2\pi\zeta,\sin2\pi\zeta,\zeta).
\]
Esta afirmación no dice que toda línea espacial de campo sea una hélice:
describe la trayectoria del extremo del vector polarización junto con el
evento marcado.

\section{Compatibilidad loxodrómica}

Para
\[
 \Phi(r,\varphi,z,t)
 =\ell\varphi+\nu\log(r/r_0)+k_zz-\omega t,
\]
la transformación
\[
 (r,\varphi)\longmapsto
 ((2+\sqrt3)r,\varphi+\pi/2)
\]
preserva la fase módulo \(2\pi\) exactamente cuando
\[
 \nu\log(2+\sqrt3)+\ell\frac{\pi}{2}\in2\pi\Z.
\]
La dilatación se interpreta como cambio de escala entre soluciones; no
como producción espontánea de energía.

\section{Centro electrónico y lector posterior}

El electrón ocupa el punto fijo único \((5,5)\) del atlas material. Esta
centralidad no se deriva otra vez en el presente trabajo. La compatibilidad
operatoria previa incluye
\[
 C^{(3)}m_{\ph}=-\frac12m_{\ph},
 \qquad
 1-\left(\frac12\right)^2=\frac34=\frac{90}{120},
\]
y
\[
 \|[P_{\Sigma,3},P_{\Pi,3}]\|=\frac{\sqrt3}{4}.
\]
La teoría de correspondencias espirales añade un lector de rutas; no
reemplaza la Ley General de Masas ni usa el electrón para seleccionar
retrospectivamente \(p\), \(\Lambda\) o \(C\).

\section{Matriz cerrada de controles negativos}

\begin{longtable}{p{.38\textwidth}p{.50\textwidth}}
\caption{Falsadores y errores de tipo excluidos por la construcción.}
\label{espiral:tab:falsadores-ap}\\
\toprule
Fallo&Consecuencia\\
\midrule
\endfirsthead
\toprule
Fallo&Consecuencia\\
\midrule
\endhead
La curva clásica selecciona \(p,\Lambda,C\).&Circularidad: el lector deja
de ser constructivo--generativo.\\
Se conserva sólo la holonomía final.&Pérdida de prefijos, carry y frontera;
no se acredita naturalidad.\\
Se exige invariancia bajo \(\Gamma_9\).&Se borra el avance de memoria; la
relación correcta es equivarianza.\\
Se confunden \(U_6\), \(w_6\) y \(\F_3^6\).&Colapso de objetos con
cardinalidades \(468\), \(243\) y \(729\).\\
Se identifica retorno de fase con retorno de estado.&Contradicción con
\(B_{K+9}\neq B_K\) en general.\\
Se identifica \(\tau\), \(p\) y \(\kappa_{\mathrm{F}}\).&Colapso de régimen,
carácter radial y curvatura visible.\\
Se llama hélice general a cualquier radio variable.&Afirmación diferencial
sin verificar la condición de Lancret.\\
Se identifica el eje \(\varphi\) con el electrón.&Mezcla de la fibra de
constantes y la realización material.\\
Se interpreta Pell como amplificación energética.&Violación de la lectura
de escala y de la conservación física.\\
Se identifica helicidad electromagnética con espín \(1/2\).&Confusión entre
representaciones distintas.\\
Se impone una celda puntual a toda especie.&Contradicción con el carácter
de fibra u multisección de las especies no centrales.\\
Se separan las cinco construcciones del continuo.&Ruptura de la genealogía
única anterior a cualquier publicación geométrica.\\
\bottomrule
\end{longtable}

\chapter{Glosario matemático}

\begin{longtable}{p{.22\textwidth}p{.66\textwidth}}
\caption{Símbolos y términos rectores.}
\label{espiral:tab:glosario}\\
\toprule
Símbolo o término&Definición\\
\midrule
\endfirsthead
\toprule
Símbolo o término&Definición\\
\midrule
\endhead
\(\Sigma,\Pi\)&Evaluaciones aditiva y multiplicativa sobre el soporte
común APP.\\
\(\tau\)&Régimen local de orientación TRIT; no es el grado radial ni la
torsión de Frenet.\\
\(M_{\ph}\)&Selector trítico de fase operativa.\\
\(\Gamma_9\)&Holonomía nonádica de la dinámica enriquecida.\\
\(H_N\)&Holonomía afín acumulada hasta el prefijo \(N\).\\
\(\Lambda_N,\mathcal C_N\)&Parte lineal y cociclo traslacional de \(H_N\).\\
\(p\)&Cociclo entero del carácter radial.\\
\(a\)&Carácter orientado local que separa genealogías del mismo
grado local.\\
\(L_p,\exp_p\)&Coordenada potencia--logarítmica y su inversa positiva.\\
\(\boxplus_p\)&Ley deformada que publica la interacción suma--producto.\\
\(\Led_N\)&Ledger enriquecido de hoja, carry, frontera, ruta y memoria.\\
\(\mathcal P_{\rad}^{\enr}\)&Lector radial enriquecido anterior a la curva.\\
\(\mathcal P_{\arq}^{\mathrm{esp}}\)&Evaluación posterior en una curva de
\(\R^2\).\\
\(\kappa_{\mathrm{F}},\mathcal T_{\mathrm{F}}\)&Curvatura y torsión de Frenet de la
curva visible.\\
Suspensión helicoidal radial&Superficie u hoja de radio variable que
conserva fase, profundidad y recorrido.\\
Publicación no fiel&Evaluación que puede identificar historias distintas.\\
Constructivo--generativo&Método que construye el objeto y su genealogía
antes de reconocer su apariencia convencional.\\
\bottomrule
\end{longtable}
